\documentclass[../DoAn.tex]{subfiles}
\begin{document}

Phụ lục này đặc tả các use case bổ sung không được trình bày đầy đủ trong Chương \ref{chapter:Related_works} do giới hạn về không gian. Các use case dưới đây thuộc nhóm quản lý thiết bị và quản lý vùng an toàn -- hai nhóm chức năng có cấu trúc CRUD tương tự nhau.

\section{Đặc tả use case: Đổi tên thiết bị}

\begin{table}[H]
\centering
\caption{Đặc tả use case Đổi tên thiết bị}
\label{table:uc_rename_device}
\begin{tabular}{|l|p{9.5cm}|}
\hline
\textbf{Tên use case}       & Đổi tên thiết bị \\
\hline
\textbf{Tác nhân}           & Phụ huynh \\
\hline
\textbf{Tiền điều kiện}     & Phụ huynh đã đăng nhập; thiết bị đã được đăng ký với tài khoản. \\
\hline
\textbf{Luồng sự kiện chính} &
1. Phụ huynh chọn thiết bị cần đổi tên trong danh sách. \newline
2. Phụ huynh chọn chức năng chỉnh sửa và nhập tên mới. \newline
3. Hệ thống cập nhật tên thiết bị trong cơ sở dữ liệu. \newline
4. Ứng dụng hiển thị tên mới ngay lập tức trên danh sách và bản đồ. \\
\hline
\textbf{Luồng phát sinh}    & Nếu tên mới để trống, hệ thống giữ nguyên tên cũ. \\
\hline
\textbf{Hậu điều kiện}      & Thiết bị hiển thị với tên mới trên toàn bộ giao diện ứng dụng. \\
\hline
\end{tabular}
\end{table}

\section{Đặc tả use case: Xóa thiết bị}

\begin{table}[H]
\centering
\caption{Đặc tả use case Xóa thiết bị}
\label{table:uc_delete_device}
\begin{tabular}{|l|p{9.5cm}|}
\hline
\textbf{Tên use case}       & Xóa thiết bị \\
\hline
\textbf{Tác nhân}           & Phụ huynh \\
\hline
\textbf{Tiền điều kiện}     & Phụ huynh đã đăng nhập; thiết bị đang được liên kết với tài khoản. \\
\hline
\textbf{Luồng sự kiện chính} &
1. Phụ huynh chọn thiết bị cần xóa trong danh sách. \newline
2. Phụ huynh chọn chức năng xóa; ứng dụng hiển thị hộp thoại xác nhận. \newline
3. Phụ huynh xác nhận xóa. \newline
4. Hệ thống hủy liên kết thiết bị khỏi tài khoản. \\
\hline
\textbf{Luồng phát sinh}    & Nếu phụ huynh hủy tại bước xác nhận, không có thay đổi nào được thực hiện. \\
\hline
\textbf{Hậu điều kiện}      & Thiết bị không còn xuất hiện trong danh sách; dữ liệu lịch sử vị trí liên quan được giữ nguyên trong cơ sở dữ liệu. \\
\hline
\end{tabular}
\end{table}

\section{Đặc tả use case: Sửa vùng an toàn}

\begin{table}[H]
\centering
\caption{Đặc tả use case Sửa vùng an toàn}
\label{table:uc_edit_geofence}
\begin{tabular}{|l|p{9.5cm}|}
\hline
\textbf{Tên use case}       & Sửa vùng an toàn \\
\hline
\textbf{Tác nhân}           & Phụ huynh \\
\hline
\textbf{Tiền điều kiện}     & Phụ huynh đã đăng nhập; tồn tại ít nhất một vùng an toàn trong tài khoản. \\
\hline
\textbf{Luồng sự kiện chính} &
1. Phụ huynh chọn vùng an toàn cần chỉnh sửa trên bản đồ hoặc trong danh sách. \newline
2. Phụ huynh thay đổi các thông số: tên, bán kính (đối với vùng tròn) hoặc các điểm tọa độ (đối với chế độ đường đi), hoặc trạng thái kích hoạt. \newline
3. Phụ huynh xác nhận lưu thay đổi. \newline
4. Hệ thống cập nhật cấu hình vùng an toàn và bắt đầu áp dụng ngay cho các bản tin vị trí tiếp theo. \\
\hline
\textbf{Luồng phát sinh}    & Nếu phụ huynh vô hiệu hóa vùng an toàn, hệ thống tạm dừng kiểm tra vi phạm cho vùng đó mà không xóa cấu hình. \\
\hline
\textbf{Hậu điều kiện}      & Vùng an toàn được cập nhật; backend áp dụng cấu hình mới từ bản tin vị trí tiếp theo nhận được. \\
\hline
\end{tabular}
\end{table}

\section{Đặc tả use case: Xóa vùng an toàn}

\begin{table}[H]
\centering
\caption{Đặc tả use case Xóa vùng an toàn}
\label{table:uc_delete_geofence}
\begin{tabular}{|l|p{9.5cm}|}
\hline
\textbf{Tên use case}       & Xóa vùng an toàn \\
\hline
\textbf{Tác nhân}           & Phụ huynh \\
\hline
\textbf{Tiền điều kiện}     & Phụ huynh đã đăng nhập; tồn tại ít nhất một vùng an toàn trong tài khoản. \\
\hline
\textbf{Luồng sự kiện chính} &
1. Phụ huynh chọn vùng an toàn cần xóa và chọn chức năng xóa. \newline
2. Ứng dụng hiển thị hộp thoại xác nhận. \newline
3. Phụ huynh xác nhận xóa. \newline
4. Hệ thống xóa cấu hình vùng an toàn; backend ngừng kiểm tra vi phạm cho vùng đó ngay lập tức. \\
\hline
\textbf{Luồng phát sinh}    & Nếu phụ huynh hủy tại bước xác nhận, không có thay đổi nào được thực hiện. \\
\hline
\textbf{Hậu điều kiện}      & Vùng an toàn không còn xuất hiện trên bản đồ; lịch sử cảnh báo liên quan được giữ nguyên. \\
\hline
\end{tabular}
\end{table}

\end{document}
